This section measures the moral hazard and selection effects associated with the monitoring program. Consumers who opt in to monitoring become safer when monitored. Despite this behavioral change, monitoring reveals persistent and previously unobserved risk differences across consumers, which in turn drive advantageous selection into the program.

\subsection{Risk Reduction and the Moral Hazard Effect\label{subsec:rf_incentive}}

If monitoring technology can effectively capture accident risk and individuals' risk is modifiable, then we should expect monitored drivers to be safer during the initial monitoring period than subsequent unmonitored ones---a (reverse) moral hazard effect.\footnote{See \textcite{Fama1980,Holmstrom1999}. A similar setting is studied in \textcite{Taylor2004, Fudenberg2006}: consumers may refrain from buying expensive items online if they know doing so will label them as inelastic shoppers and lead to higher future prices.}

We now directly measure this effect by comparing claim outcomes for monitored consumers before and after monitoring ends.
We construct a balanced panel over the first three periods (18 months) consisting of all consumers that are eligible for monitoring, and adopt a \textit{difference-in-differences} approach to control for spurious trends in claim rates across periods that are irrelevant to monitoring. Among monitored drivers, we take the first difference in claim counts\footnote{Throughout our reduced-form analyses, we use claim count as our cost proxy. This is because claim severity is extremely noisy and skewed. This is also common practice in the industry, where many risk-rating algorithms are set to predict risk occurrence only. We therefore present our estimates mostly in percentage comparison terms.} between the monitoring and post-monitoring periods. This difference is then benchmarked against its counterpart among unmonitored consumers (the control group). Our main specification adds exhaustive observable controls and coverage fixed effects:
\begin{align}
C_{it}= & \alpha+\tau m_{i}+\omega\mathbf{1}_{post,t}+\theta_{mh}m_{i}\cdot\mathbf{1}_{post,t}+\mathbf{x}_{it}^{\prime}\mathbf{\beta}+\mathbf{y}_{it}^{\prime}\mathbf{\psi}+\epsilon_{it}\label{eq:mh_reg}
\end{align}
Here, $i$ and $t$ index the consumer and period in our panel dataset, respectively. $C$ denotes the claim count, and $m_i$ is a consumer-specific indicator for whether $i$ finished monitoring. The vector $\mathbf{x}$ denotes a rich set of observable characteristics that the firm uses in pricing, which includes state fixed effects and variables outlined in Table \ref{tab:sum_stat_x}. $y$ denotes fixed effects for the liability limit chosen as well as an indicator for whether the consumer has purchased property coverage at the start of each period, as summarized in Table \ref{tab:sum_stat}.
Last, we add a specification controlling for driver fixed effects to account for selection on consumer unobservables.

\inputifexists{\exhibitpath/tab_2.tex}

Among consumers who are monitored, the actual duration of monitoring varies due to logistical discrepancies such as mailing and installation times; some drivers also start the program but drop out before finishing.
We thus introduce an additional specification that adds treatment intensity $z_{i}$ and the interaction with monitoring timing to our main specification, using all drivers who are monitored. Here, $z_{i}$ is calculated as the fraction of days monitored in the first period minus the same fraction in post-monitoring periods. \footnote{Monitoring duration may be correlated with claims due to reasons other than moral hazard, and particularly if claims itself leads to a non-finishing status. However, only 13 claims occurred within 7 days before or after monitoring drop-out out of more than a hundred claims that we observe for monitored drivers who did not finish. Further, if some monitored consumers drop out as they discover that they cannot change their risk, our estimate here would encompass both the moral hazard effect and a selection on moral hazard effect \textit{a la} \textcite{Einav2013}.}

Results are reported in
\Cref{tab:mh_results}. We find a large moral hazard effect. Column 3 corresponds to the specification in \Cref{eq:mh_reg}, which shows that the average claim count for monitored consumers is 0.007 or 21.93\% higher after the monitoring period. Adjusting for the average monitoring duration of first-period monitoring finishers of 131 days---only a fraction of the monitoring period---the moral hazard effect would be 29.97\%. This result is stable across specifications and only gets slightly bigger as we add controls; for instance, adding additional variation in monitoring duration (treatment intensity) generates a similar result of 33.08\% (derived by setting $z=1$ and adding up the two interaction coefficients in column 7). We test for parallel trends between the monitored and unmonitored groups by repeating our main specification in balanced panels encompassing subsequent unmonitored periods, which also serves as a placebo check. As columns 9-11 show, no differential claim change across periods can be detected.

Our main results above use a balanced panel covering the first three periods of drivers' engagement with the monitoring firm only. Below, we double the time horizon to six periods (three years) and use the full unbalanced panel. We adapt specification \ref{eq:mh_reg} slightly to look at the claim progression over each of the six periods across monitoring groups. The fixed-effect estimates for each policy period, $\theta_t$, are illustrated in \Cref{fig:rf-illus}. The level difference between the grey and the orange lines after period 0 represents a persistent difference in risk between opt-in consumers and their opt-out counterparts---a selection effect quantified in the next subsection. There is an additional risk reduction in period 0 among opt-in consumers, which closely tracks the moral hazard effect shown in Table \ref{tab:mh_results} column (3).
\begin{align}
C_{it}= & \alpha+\tau m_{i}+\omega_t\mathbf{1}_{t}+\theta_{t}m_{i}\cdot\mathbf{1}_{t}+\mathbf{x}_{it}^{\prime}\mathbf{\beta}+\epsilon_{it}\label{eq:mh_reg_long}
\end{align}
\vspace{-0.5cm}
\begin{figure}[htbp!]
\centering
\includeifexists[scale=0.5]{\exhibitpath/fig_4.png}
\caption{Claim Progression across Monitoring Groups \label{fig:rf-illus}}
\vspace{0.1cm}
\begin{minipage}{0.95\textwidth} % choose width suitably
    {\footnotesize \emph{Notes: }This graph reports the fixed effect estimates of \cref{eq:mh_reg_long} using the full unbalanced panel. In Appendix \ref{subsec:app-unbal}, we show that our results are robust to using a six-period balanced panel. The grey line plots $\omega_t$ while the orange line plots $\omega_t + \theta_t$, both against insurance periods $t$. The red box is superimposed ex-post to represent the period when opt-in consumers are monitored. Error-bars report 95\% confidence interval. \par}
\end{minipage}
\par
\end{figure}
\vspace{-1em}

\paragraph{Robustness and Discussion}
We show in Appendix \ref{subsec:app-unbal} that our results in this section are robust to whether we use a balanced panel or not. We also investigate heterogeneity in the moral hazard effect in Appendix \ref{subsec:app-mhhet}. Overall, systematic heterogeneity across driver characteristics or plan options is rare, although we find that drivers who are initially assigned a higher risk class tend to experience a larger moral hazard effect.

We discuss two important caveats of our results. First, monitoring mitigates moral hazard because it builds a reputation mechanism by signaling consumers' persistent risk level, \textit{not} because it directly rewards effort during monitoring
. The magnitude of risk reduction can be different in the latter setting.\footnote{We are also unable to disentangle the ``Hawthorne effect'' from consumers' responsiveness to financial incentives in our estimate. Since consumers must be aware of the data collection to be incentivized for it, we consider this effect as part of the moral hazard effect.}
Second, our estimates measure a treatment-on-treated effect. Although we find little observed heterogeneity in the moral hazard effect among opt-in consumers
, it is possible that the majority of unmonitored drivers are less capable of altering their risk. The moral hazard effect that we have identified is thus likely larger than the population average.\footnote{We also suppress selection on moral hazard in counterfactuals \parencite{Einav2013}. In equilibrium, the firm assesses the signal that monitored consumers send based on their future claim records when they are no longer monitored, which corresponds to the renewal discount it gives. Thus, risk reduction is compensated only to the extent that it correlates with consumers' unmonitored risk type. If safer consumers' risk levels are also more responsive to incentives, as suggested by a pure effort cost model, selection on the incentive effect can be important. In particular, perfect revelation of a continuum of risk types is possible, as characterized in \textcite{mailath1987incentive}, with a monotonicity condition similar to the single-crossing condition. However, consumers likely have multidimensional heterogeneity in reality, so consumers' performance during monitoring may not perfectly reveal their risk types \parencite{frankel2014muddled}.} To avoid external validity concerns, our counterfactual analysis maintains the opt-in structure of the monitoring program and does not extrapolate to scenarios where the monitoring rate is significantly higher than in the data.

\subsection{Private Risk and the Selection Effect\label{subsec:rf_selection}}

Our analysis above suggests that monitored drivers remain safer than their counterparts even after the monitoring period.\footnote{For instance, the coefficients on the monitoring indicator (or monitoring intensity) in Table \ref{tab:mh_results} are large and negative, and the claim count differentials for monitored drivers, depicted in orange in Figure \ref{fig:rf-illus}, are persistently lower than their unmonitored counterparts in grey.} This suggests that the monitoring pool may be advantageously selected.  Table \ref{tab:mon_unmon_t0} reports the results of regressing claim counts in the first period only ($t=0$) on a monitoring indicator, with the same control specifications---and produces nearly identical coefficients---as columns (1) to (3) of Table \ref{tab:mh_results}. The coefficient on the monitoring indicator suggests that the risk differential between the two groups is larger than the moral hazard effect discussed in the previous section: the latter only accounts for 61.10\% of the risk differential. This suggests that consumers possess private information about their accident risk, driving advantageous selection into monitoring.
\inputifexists{\exhibitpath/tab_3.tex}

Selection into monitoring also implies that the technology is effective at capturing previously unobserved differences in drivers' risk types. In Appendix Figures \ref{fig:rf-info-part} and \ref{fig:rf-info-score}, we demonstrate that both a driver's choice to be monitored and their first-period monitoring score strongly and persistently predict future-period claim counts. For example, receiving a score that is one standard deviation above the mean is associated with a $29\%$ higher average claim count in the second period. Controlling for claims has little impact on these estimates, which suggests that the high-frequency monitoring data is highly informative of driver risk relative to sparse claims records.

A potential concern is that monitored drivers may learn to become persistently safer due to monitoring. This does not influence our moral hazard estimates, which are derived from the observed behavioral change between the monitoring period the periods afterward regardless of the pre-monitoring level of consumer risk. However, if consumers' post-monitoring risk is reduced persistently due to monitoring, ignoring such effect may exaggerate the degree of advantageous selection. In Appendix \ref{subsec:app-learning}, we use consumers' public ``at fault accident'' (AFA) records to test for this possibility. Although the AFA records are substantially noisier than our proprietary claim records, we find persistent gaps between the average risk captured by AFAs among monitored and unmonitored drivers before and after they engage with the monitoring firm, consistent with our main results.

\subsection{Renewal Elasticity\label{subsec:rf_demand}}

The welfare benefit of monitoring may be limited if the firm extracts large rents during renewal periods. However, when pricing the monitoring score, the firm must consider that a monitored consumer can choose not to renew and leave for another firm. We quantify the price elasticity of renewal with the following regression:
\begin{align}
    \mathbf{1}^{\text{renewed}}_{it} & = \alpha M_{i} + \theta \log p^{\text{renewal}}_{it} + \theta_M M_{i} \cdot \log p^{\text{renewal}}_{it} + \psi s_{i} + \mathbf{x}_{it}^\prime\beta + \epsilon_{it}.\label{eq:rf_demand_2S}
\end{align}
The dependent variable indicates whether consumer $i$ renewed their insurance plan in period $t$. We augment the consumer-specific monitoring indicator $m$ to get $M$, which admits three discrete monitoring \textit{groups}: unmonitored consumers, monitored consumers that received discounts, and monitored consumers who received no discount or even a surcharges. We include the same controls $\mathbf{x}$ as Column (3) of Table \ref{tab:mh_results}. We also add the consumer-specific monitoring score $s_i$, normalized among monitored consumers and set to $0$ for unmonitored ones. The parameters of interest are $\theta$ and $\theta_M$, which combines with $M$ to give price elasticities across monitoring groups.

The regression above may produce biased estimates if renewal prices are endogenous to unobserved demand factors or competitor pricing. For robustness, we adopt an instrumental-variable approach to compute price elasticity. We narrow the sample down to 30-day windows around rate revision events within each state, and instrument for renewal pricing using $Z_{it}$, an indicator for whether the consumer arrived after the rate revision in their market took effect. Our estimation adds rate revision fixed effects to the controls of both the first-stage and the reduced-form regressions (Equations \ref{eq:rf_demand_2S} and \ref{eq:rf_demand_1S}). The exclusion restriction is supported by an event-study argument: consumers arriving right before and after a rate revision event should be similar, and the price differential they face is thus independent from their idiosyncratic unobservable demand factors.
\begin{align}
    \log p^{\text{renewal}}_{it} & = \tau M_{i} + \omega Z_{it} + \omega_M M_{i} \cdot Z_{it} + \phi s_{i} + \mathbf{x}_{it}^\prime \gamma + \nu_{it}.
    \label{eq:rf_demand_1S}
\end{align}

\Cref{fig:rf-demand} shows that both OLS and IV regressions produce similar price elasticity estimates across monitoring groups.
During renewal, monitored consumers who receive discounts are less price sensitive than the average unmonitored consumer, while the opposite is true for those who receive no discount or who receive a surcharge. This implies that the firm faces a steeper residual demand curve when giving discounts and a flatter one when imposing surcharges. For safe monitored drivers---those who proved to be less risky than the firm initially expected---this suggests that the firm can extract an informational rent by raising markups after the monitoring period.
However, it is unclear how the firm should react to the higher price sensitivity among monitored but discount-ineligible consumers. While steering the riskiest drivers toward competitors is likely profitable, many risky monitored drivers may still yield positive marginal profits. Moreover, the firm's ability to surcharge is directly constrained by regulatory price controls---a feature we discuss in detail in \Cref{subsec:equi_concept} after introducing our pricing model.

\begin{figure}[htbp!]
\centering
\includeifexists[scale=0.79]{\exhibitpath/fig_5.png}
\caption{Price Elasticity of Renewal Acceptance by Monitoring Groups \label{fig:rf-demand}}
\vspace{0.1cm}
\begin{minipage}{0.95\textwidth} % choose width suitably
    {\footnotesize \emph{Notes: }The dots report the renewal price elasticity estimates across monitoring groups using OLS or IV, which correspond to $\hat{\theta} + \hat{\theta_M} \cdot M$ in \Cref{eq:rf_demand_2S}. They measure the average percentage point increase in renewal likelihood with respect to 1\% increase in renewal price. $M$ are indicators for three monitoring groups (from left to right): monitored consumers that receive price discounts, monitored consumers that receive no price discounts or surcharges, unmonitored consumers. Lines reports 95\% confidence interval, with standard errors clustered on the consumer level. \par}
\end{minipage}
\end{figure}

\paragraph{The need for structural models}
Taken together, our results in this section show that monitoring produces highly informative signals on consumers' future accident risk, which effectively separates drivers who are otherwise indistinguishable to the firm. This leads to advantageous selection into the program and safer driving when consumers are monitored. We have thus separately identified the effects of moral hazard and selection.

However, in order to evaluate the welfare impact of monitoring and the trade-offs associated with data or pricing regulations, a structural model is needed for three reasons.
First, monitoring affects consumer welfare through several interacting channels that cannot be combined without estimating risk preference. Beyond an immediate reduction in accident risk,
monitored drivers may purchase higher coverage in anticipation of future discounts, but they may also be exposed to greater premium volatility due to reclassification risk \parencite{Handel2015}. Second, consumers make three discrete choices: insurance coverage, monitoring opt-in, and attrition during renewals. Their inter-dependence and correlation with consumer risk jointly determine the degree of asymmetric information (between consumers and firms and across firms), shaping pricing and efficiency in equilibrium.
Third, the pricing of monitoring is inherently dynamic and multidimensional; it is also subject to intense regulatory pressure that may have lead to overly conservative pricing.\footnote{For example, as shown in Appendix \ref{sec:app_firm_pricing}, the firm did not raise price for the unmonitored pool, which is primarily to appease some state regulators' aversion to price increase \parencite{ben2023privacy}.} Simulating how prices would change in counterfactuals therefore requires modeling the firm's main pricing levers---how it encourages consumer opt-in ex ante (to obtaining the monitoring scores), and how it limits attrition while extracting rents ex post.

In Sections \ref{sec:demand_model} and \ref{sec:estimation}, we estimate a structural model of accident risk and consumer choice. In Section \ref{sec:firm_equi}, we present a model of firm pricing across monitoring groups both pre- and post-monitoring, which facilitates welfare and counterfactual calculations.
